Un rayon lumineux passe de l'air ($n = \num{1.00}$) à un bloc transparent dont l'indice de réfraction est de \num{\v{n}}. L'angle d'incidence est de \qty{\v{th1}}{\degree}. La vitesse de la lumière dans le vide est de \qty{3.00e8}{m/s}.

\begin{center}
\begin{tikzpicture}[thick]
    \fill[gray!20] (-3,-2) rectangle (3,0);
    \draw (-3,0) -- (3,0);
    \draw[thin,dashed] (0,2.3) -- (0,-2);
    \draw[-{Stealth[length=3mm]}] ({90+\v{th1}}:2.8) -- ({90+\v{th1}}:1.2);
    \draw ({90+\v{th1}}:1.2) -- (0,0);
    \draw[thin] (0,1.3) arc (90:{90+\v{th1}}:1.3);
    \path (0,0) ++({90+\v{th1}/2}:1.75) node {\small\ang{\v{th1}}};
    \node[right] at (0.8,1.2) {\small air ($n = \num{1.00}$)};
    \node[right] at (0.8,-1) {\small bloc ($n = \num{\v{n}}$)};
\end{tikzpicture}
\end{center}

\begin{enumerate}[label=\alph*)]
    \item Déterminez l'angle de réfraction dans le bloc.
    \item Déterminez la vitesse de la lumière dans le bloc.
    \item De combien de degrés le rayon a-t-il été dévié ? Le rayon s'est-il rapproché ou éloigné de la normale ?
\end{enumerate}
